\section{Discussion}
\label{sec:discussion}

Reliable maintenance of a relative near-optimal set imposes a cost that is
not visible from allocation regret alone. One arrival can affect many reports
through uncertainty about the common reference level. Repeated improvements
can also make fine discrimination of earlier candidates unnecessary. The
cumulative area criterion measures both effects in the same experiment
and states how much unresolved time a reporting requirement permits.

The strength of the timeliness requirement matters. Sublinear area here means
that the average unresolved \emph{count} vanishes. An application concerned
only with the unresolved proportion, the discovery of one good arm, or a
terminal decision would face a different task. The arrival exponent and
quadratic worst-case area depend on this normalization. Arms in the boundary
band remain in the model: their unresolved time is exempt from the loss, while
every issued label remains subject to correctness.

The finite-time results and eventual consistency also answer different
questions. If arrivals stop, the final arm set is finite, and every arm is
strictly separated from the membership boundary, a revisable multiresolution
version eventually reports the exact set almost surely and has sublinear
expected regret. Its construction and proof are supplied in the supplement.
For the fixed-precision checkpoint implementation, frozen intervals can remain
too wide and an early reporting error need not be corrected, so this consistency
conclusion does not apply. The reward layer retains every available arm and
can still correct allocation.

Our bounds leave a substantive instancewise question. Under arbitrary
overlapping arrivals, heterogeneous boundary distances, and intermediate area
budgets, current information constraints need not describe an attainable causal
allocation. A matching theory would have to price common reference evidence,
unfinished tasks, their possible expiration, and observations already obtained
for reward collection without charging the same observation twice. The
restricted multiple-update theorem leaves that scheduling problem open.
The critical logarithmic gap on an individual smooth schedule and automatic
choice of service precision also remain open.

For practice, the numerical evidence favors examining the two costs together
and comparing methods that retain their history. CP-UCB-AA gives transparent
guarantees, but fixed checkpoints can devote many observations to tasks already
unimportant at a moderate horizon. Confidence-driven adaptations can be much
more efficient, motivating sharper continual-service guarantees for those
rules. A smaller reporting set alone does not establish a resource advantage.

\paragraph{Supplementary material and reproducibility.}
The supplement contains complete proofs, the information-allocation lower
relaxation, the quantile extension, the revisable consistency construction,
and further allocation diagnostics. The numerical package contains the policy
definitions, frozen configurations, every replication, paired Monte Carlo
summaries, figure code, and a source-class parity check. All numerical data in
this article are simulated; no patient-level or other original human-subject
data were used.
